\documentclass[../main.tex]{subfiles}

\begin{document}
\subsection{Funciones armónicas}
\classdate{2026-04-08}

\begin{definition}
	Decimos que \(f \colon U \subseteq \mathbb{R}^{2} \to \mathbb{R}\) es armónica
	\(\bigl(f \in C^{2}\bigr)\) si

	\begin{equation}
		\Delta{f} = \pdv[order = 2]{f}{x} + \pdv[order = 2]{f}{y} = 0.
	\end{equation}
\end{definition}

\begin{proposition}
	Los siguientes enunciados son equivalentes para \(f \colon U \subseteq \mathbb{C}
	\to \mathbb{C}\) con \(f = u + iv\),

	\begin{enumerate}
		\item \(f\) es holomorfa.
		\item \(u,v\) son armónicas si cumplen Cauchy-Riemann.
	\end{enumerate}
\end{proposition}

\begin{proof}
	\begin{itemize}
		\item (a) \(\,\implies\,\) (b)
		      Como \( f\) es holomorfa, \(u\) y \(v\) cumplen Cauchy-Riemann y además son
		      \(C^{2}\).
		      Ahora notemos que

		      \begin{align*}
			      \Delta{u} = \pdv[order = 2]{u}{x} + \pdv[order = 2]{u}{y} & =
			      \pdv*{\pdv{u}{x}}{x} + \pdv*{\pdv{u}{y}}{y},                                                                        \\
			                                                                & = \pdv*{\pdv{v}{y}}{x} - \pdv*{\pdv{v}{x}}{y},\nonumber \\
			                                                                & = \pdv{v}{x,y} - \pdv{v}{x,y} = 0.
		      \end{align*}

		      Y \(\Delta{v} = 0\) es análogo. Por lo tanto, \(u,v\) son armónicas y cumplen
		      Cauchy-Riemann.

		\item (b) \(\,\implies\,\) (b)

		      Que \(u\) y \(v\) sean armónicas implica que son \(C^{2}\) y por tanto \(f\) vista como
		      función \(\mathbb{R}^{2} \to \mathbb{R}^{2}\) es derivable y además como cumple
		      Cauchy-Riemann, por un teorema visto en clase, \(f\) es holomorfa.
	\end{itemize}
\end{proof}

\begin{definition}
	Decimos que dos funciones armónicas \(u,v \colon U \subseteq \mathbb{R} \to
	\mathbb{R}\) son conjugadas armónicas una de la otra si cumplen Cauchy-Riemann. O
	bien, que la función \( f \colon U \subseteq \mathbb{C} \to \mathbb{C}\) dada por
	\(u + iv\) sea analítica.
\end{definition}

\begin{example}
	Sean \(u,v \colon \mathbb{C} \to \mathbb{R}\) con \(u(x + iy) = x^{2} - y^{2}\) y
	\(v(x + iy) = 2xy\). Entonces,

	\begin{align*}
		\pdv{u(x, y)}{x} & = \pdv{v(x,y)}{y},   \\
		\pdv{u(x, y)}{y} & = - \pdv{v(x,y)}{x}.
	\end{align*}

	Por lo tanto cumplen Cauchy-Riemann y son conjugadas armónicas una de la otra.
\end{example}

Sea \(u \colon B_{r}(z_{0}) \subseteq \mathbb{C} \to \mathbb{R} \subseteq \mathbb{C},
z_{0} = u_{0} + iv_{0}\) es armónica y queremos encontrar otra función armónica \(v
\colon B_{r}(z_{0}) \to \mathbb{R}\) tal que cumpla Cauchy-Riemann.
Supongamos \(v\), \(\pdv{v(x,y)}{y} = \pdv{u(x, y)}{x}\),

\begin{equation*}
	v(x, y) = \int_{y_{0}}^{y}\odif[sep-end=\medspace]{dt} \pdv{u(x, t)}{x} + h(x).
\end{equation*}

Entonces,

\begin{align*}
	-\pdv{u(x, y)}{y} = \pdv{v(x, y)}{x} & =
	\pdv{}{x}\int_{y_{0}}^{y}\odif[sep-end=\medspace]{t} \pdv{u(x, t)}{x} +
	h^{\prime}(x),                                                                       \\
	                                     & = \int_{y_{0}}^{y}\odif[sep-end=\medspace]{t}
	\pdv[order = 2]{u(x, t)}{x} +
	h^{\prime}(x),\shortintertext{pues \(u\) es armónica}
	                                     & = -
	\int_{y_{0}}^{y}\odif[sep-end=\medspace]{t}
	\pdv[order = 2]{u(x, t)}{y} + h^{\prime}(x),                                         \\
	                                     & = - \pdv{u(x,y)}{y} + \pdv{u(x, y_{0})}{y} +
	h^{\prime}(x).
\end{align*}

Por lo que

\begin{align*}
	h^{\prime}(x) & = - \pdv{u(x, y_{0})}{y},                                                 \\
	h(x)          & = - \int_{x_{0}}^{x}\odif[sep-end=\medspace]{s} \pdv{u(s, y_{0})}{y} + C.
\end{align*}

Y por lo tanto,

\begin{equation*}
	v(x, y) = \int_{y_{0}}^{y}\odif[sep-end=\medspace]{t} \pdv{u(x, t)}{x} -
	\int_{x_{0}}^{x}\odif[sep-end=\medspace]{s} \pdv{u(s, y_{0})}{y} + C.
\end{equation*}
\end{document}
